% Logros del laboratorio 3 frente a sus objetivos y trabajo del docente.
% Generado por ../generar_logros.py desde ../logros_datos.py: editar los datos y regenerar.
\documentclass[a4paper,11pt]{article}
% Preámbulo común de los informes de laboratorio (II/2026). Se incluye desde laboratorios/labNN_*/
%   % Preámbulo común de los informes de laboratorio (II/2026). Se incluye desde laboratorios/labNN_*/
%   % Preámbulo común de los informes de laboratorio (II/2026). Se incluye desde laboratorios/labNN_*/
%   \input{../informe_comun}
\input{../../gemelo-digital/manual/estilo_manual}
\input{../../gemelo-digital/manual/tecnico/macros}
\graphicspath{{../../documentacion/planos_isa/pid/}{../../documentacion/planos_isa/electrico/}{../../documentacion/manuales_inversa/}
  {../../database/mockups/}{../../database/documentacion/}{../../database/modelo/}{../../database/interfaz/}
  {../../gemelo-digital/manual/img/}{../../modelo-3d/informe/figuras/}{../../modelo-3d/levantamiento/}{../../../../global/fotos-planta/}{../figuras/}}
\usepackage{pdflscape}
\lstset{showstringspaces=false}
% plano A3 en una página apaisada: \plano[página]{archivo}{leyenda}
\newcommand{\plano}[3][1]{\begin{landscape}\begin{figure}[p]\centering\includegraphics[page=#1, width=\linewidth, height=0.86\textwidth, keepaspectratio]{#2}\caption{#3}\end{figure}\end{landscape}}
\setlength{\headheight}{14pt}
\newcommand{\nlab}{}
\newcommand{\tlab}{}
\lhead{\small Brazo XK-JS1-1 · Informe del laboratorio \nlab}
\rhead{\small UMSA · II/2026}

% \portada{número}{título}{inicio}{entrega}{semanas}
\newcommand{\portada}[5]{%
\renewcommand{\nlab}{#1}%
\begin{titlepage}
\centering
\vspace*{0.6cm}
{\large Universidad Mayor de San Andrés · Facultad de Ingeniería\par}
{\large Carrera de Ingeniería Electrónica\par}
\vspace{0.8cm}
{\color{naranja!80!black}\rule{\linewidth}{1.5pt}\par}
\vspace{0.3cm}
{\Large Informe del Laboratorio Integrado Nº #1\par}
\vspace{0.2cm}
{\huge\bfseries #2\par}
\vspace{0.3cm}
{\color{naranja!80!black}\rule{\linewidth}{1.5pt}\par}
\vspace{0.8cm}
{\large Planta: brazo robótico XK-JS1-1 de 4 grados de libertad\par}
\vspace{0.3cm}
{\large ETN-902 Sistemas de Control II · ETN-1034 Aplicación de Técnicas de Control · ETN-1000 Bases de Datos\par}
\vspace{1cm}
\begin{tabular}{rl}
\textbf{Docente:} & Ing. Jorge Antonio Nava Amador\\
\textbf{Auxiliares:} & Limber Ajnota · Fabián Plata · Franz Choque\\
\textbf{Inicio de trabajo:} & #3\\
\textbf{Entrega:} & #4 (#5)\\
\end{tabular}
\vspace{1cm}

\begin{tabular}{ll}
\multicolumn{2}{l}{\textbf{Escuadra del brazo de 4 GDL}}\\[2pt]
Control (ETN-902 / ETN-1034): & Gabriel Andres Gomez Calla\\
 & Jose Lenin Mamani Choque\\
 & Miguel Alejandro Lima Huarachi\\
 & Román Abner Machicado Arequipa\\[2pt]
Cómputo (ETN-1000): & Miguel Alejandro Lima Huarachi\\
 & Marco Eduardo Chura Rondo\\
\end{tabular}
\vfill
{\small Repositorio: \url{https://github.com/GabAndGom/Brazo-de-4-grados-de-libertad}\par}
{\small La Paz, Bolivia · gestión II/2026\par}
\end{titlepage}
}

% Tabla de entregables (Definition of Done): \begin{dod} ... \end{dod}, filas: entregable & estado & evidencia\\
\newenvironment{dod}{\begin{center}\small\begin{longtable}{>{\raggedright}p{5.2cm} >{\raggedright}p{1.8cm} >{\raggedright\arraybackslash}p{7.6cm}}
\toprule \textbf{Entregable} & \textbf{Estado} & \textbf{Evidencia / observación}\\ \midrule \endhead}
{\bottomrule\end{longtable}\end{center}}

\input{../../gemelo-digital/manual/estilo_manual}
% Macros del manual técnico: ladder, estados, código y rutas de imágenes
\usetikzlibrary{shapes.multipart}
\graphicspath{{../img/}{../../../../../global/fotos-planta/}{../../../modelo-3d/informe/figuras/}{../../../modelo-3d/levantamiento/}{../../../database/modelo/}{../../../documentacion/manuales_inversa/}{../../../documentacion/planos_isa/electrico/}{../../../database/interfaz/}}

\lstdefinestyle{py}{language=Python, keywordstyle=\color{azulg}\bfseries, commentstyle=\color{grisos}\itshape,
  stringstyle=\color{verde}, basicstyle=\ttfamily\scriptsize, numbers=left, numberstyle=\tiny\color{grisos}}
\lstdefinestyle{js}{language=Java, morekeywords={const,let,function,async,await,return,for,of,if,else}, keywordstyle=\color{azulg}\bfseries,
  commentstyle=\color{grisos}\itshape, stringstyle=\color{verde}, basicstyle=\ttfamily\scriptsize, numbers=left, numberstyle=\tiny\color{grisos}}
\lstdefinestyle{awl}{basicstyle=\ttfamily\tiny, commentstyle=\color{grisos}\itshape, morecomment=[l]{//},
  keywordstyle=\color{azulg}\bfseries, morekeywords={LD,LDN,A,AN,O,ON,EU,ED,S,R,MOVB,MOVW,MOVD,PLS,TON,CALL,CRET,LPS,LRD,LPP,ALD,OLD,FILL,LDW,LDB,LDD,AW,AD,AB,OW,OD},
  numbers=left, numberstyle=\tiny\color{grisos}, frame=single, backgroundcolor=\color{gris}}
\lstdefinestyle{sql}{language=SQL, keywordstyle=\color{azulg}\bfseries, commentstyle=\color{grisos}\itshape,
  stringstyle=\color{verde}, basicstyle=\ttfamily\tiny, numbers=left, numberstyle=\tiny\color{grisos},
  morekeywords={REFERENCES,SERIAL,SMALLSERIAL,BIGSERIAL,TIMESTAMPTZ,BOOLEAN,INET,RETURNS,TRIGGER,LANGUAGE,DECLARE,RAISE,EXCEPTION,VIEW,INDEX}}

\tikzset{
  ld/.style={line width=0.7pt},
  pics/na/.style={code={\fill[white](-0.12,-0.26) rectangle (0.12,0.26);
     \draw[thick](-0.1,-0.24)--(-0.1,0.24)(0.1,-0.24)--(0.1,0.24);
     \node[font=\scriptsize\ttfamily, above] at (0,0.22){#1};}},
  pics/nc/.style={code={\fill[white](-0.12,-0.26) rectangle (0.12,0.26);
     \draw[thick](-0.1,-0.24)--(-0.1,0.24)(0.1,-0.24)--(0.1,0.24);\draw(-0.17,-0.24)--(0.17,0.24);
     \node[font=\scriptsize\ttfamily, above] at (0,0.22){#1};}},
  pics/pf/.style={code={\fill[white](-0.12,-0.26) rectangle (0.12,0.26);
     \draw[thick](-0.1,-0.24)--(-0.1,0.24)(0.1,-0.24)--(0.1,0.24);
     \node[font=\tiny\bfseries] at (0,0){P};}},
  pics/cmp/.style n args={3}{code={\fill[white](-0.12,-0.26) rectangle (0.12,0.26);
     \draw[thick](-0.1,-0.24)--(-0.1,0.24)(0.1,-0.24)--(0.1,0.24);
     \node[font=\tiny\ttfamily] at (0,0){#2};
     \node[font=\scriptsize\ttfamily, above] at (0,0.22){#1};
     \node[font=\scriptsize\ttfamily, below] at (0,-0.22){#3};}},
  pics/bob/.style n args={3}{code={\fill[white](-0.2,-0.26) rectangle (0.2,0.26);
     \draw[thick](-0.05,0.24) arc (110:250:0.255);\draw[thick](0.05,0.24) arc (70:-70:0.255);
     \node[font=\tiny\bfseries] at (0,0){#2};
     \node[font=\scriptsize\ttfamily, above] at (0,0.22){#1};
     \node[font=\scriptsize\ttfamily, below] at (0,-0.22){#3};}},
  fb/.style={draw, fill=white, font=\scriptsize\ttfamily, align=left, inner sep=3pt, anchor=west},
  barra/.style={line width=2pt},
  estado/.style={draw, rounded corners=3pt, fill=white, align=center, font=\scriptsize, minimum width=22mm, minimum height=10mm, inner sep=3pt},
  estadon/.style={estado, fill=naranja!14, draw=naranja!70!black},
  trans/.style={font=\scriptsize, align=center, fill=white, inner sep=1pt},
  lugar/.style={circle, draw, thick, minimum size=7mm, fill=white},
  transicion/.style={rectangle, fill=black, minimum width=1.6mm, minimum height=7mm, inner sep=0},
  tabla/.style={rectangle split, rectangle split parts=2, draw, fill=white, font=\tiny\ttfamily, align=left, inner sep=2pt, rounded corners=1pt, text width=33mm},
}
\newcommand{\pendiente}[1]{\par\smallskip\noindent\colorbox{naranja!12}{\parbox{\dimexpr\linewidth-2\fboxsep}{\textbf{\color{naranja!80!black}Pendiente.} #1}}\par\smallskip}
\newcommand{\cumple}{\textcolor{verde}{\textbf{cumplido}}}
\newcommand{\parcial}{\textcolor{naranja}{\textbf{parcial}}}
\newcommand{\falta}{\textcolor{rojo}{\textbf{pendiente}}}

\graphicspath{{../../documentacion/planos_isa/pid/}{../../documentacion/planos_isa/electrico/}{../../documentacion/manuales_inversa/}
  {../../database/mockups/}{../../database/documentacion/}{../../database/modelo/}{../../database/interfaz/}
  {../../gemelo-digital/manual/img/}{../../modelo-3d/informe/figuras/}{../../modelo-3d/levantamiento/}{../../../../global/fotos-planta/}{../figuras/}}
\usepackage{pdflscape}
\lstset{showstringspaces=false}
% plano A3 en una página apaisada: \plano[página]{archivo}{leyenda}
\newcommand{\plano}[3][1]{\begin{landscape}\begin{figure}[p]\centering\includegraphics[page=#1, width=\linewidth, height=0.86\textwidth, keepaspectratio]{#2}\caption{#3}\end{figure}\end{landscape}}
\setlength{\headheight}{14pt}
\newcommand{\nlab}{}
\newcommand{\tlab}{}
\lhead{\small Brazo XK-JS1-1 · Informe del laboratorio \nlab}
\rhead{\small UMSA · II/2026}

% \portada{número}{título}{inicio}{entrega}{semanas}
\newcommand{\portada}[5]{%
\renewcommand{\nlab}{#1}%
\begin{titlepage}
\centering
\vspace*{0.6cm}
{\large Universidad Mayor de San Andrés · Facultad de Ingeniería\par}
{\large Carrera de Ingeniería Electrónica\par}
\vspace{0.8cm}
{\color{naranja!80!black}\rule{\linewidth}{1.5pt}\par}
\vspace{0.3cm}
{\Large Informe del Laboratorio Integrado Nº #1\par}
\vspace{0.2cm}
{\huge\bfseries #2\par}
\vspace{0.3cm}
{\color{naranja!80!black}\rule{\linewidth}{1.5pt}\par}
\vspace{0.8cm}
{\large Planta: brazo robótico XK-JS1-1 de 4 grados de libertad\par}
\vspace{0.3cm}
{\large ETN-902 Sistemas de Control II · ETN-1034 Aplicación de Técnicas de Control · ETN-1000 Bases de Datos\par}
\vspace{1cm}
\begin{tabular}{rl}
\textbf{Docente:} & Ing. Jorge Antonio Nava Amador\\
\textbf{Auxiliares:} & Limber Ajnota · Fabián Plata · Franz Choque\\
\textbf{Inicio de trabajo:} & #3\\
\textbf{Entrega:} & #4 (#5)\\
\end{tabular}
\vspace{1cm}

\begin{tabular}{ll}
\multicolumn{2}{l}{\textbf{Escuadra del brazo de 4 GDL}}\\[2pt]
Control (ETN-902 / ETN-1034): & Gabriel Andres Gomez Calla\\
 & Jose Lenin Mamani Choque\\
 & Miguel Alejandro Lima Huarachi\\
 & Román Abner Machicado Arequipa\\[2pt]
Cómputo (ETN-1000): & Miguel Alejandro Lima Huarachi\\
 & Marco Eduardo Chura Rondo\\
\end{tabular}
\vfill
{\small Repositorio: \url{https://github.com/GabAndGom/Brazo-de-4-grados-de-libertad}\par}
{\small La Paz, Bolivia · gestión II/2026\par}
\end{titlepage}
}

% Tabla de entregables (Definition of Done): \begin{dod} ... \end{dod}, filas: entregable & estado & evidencia\\
\newenvironment{dod}{\begin{center}\small\begin{longtable}{>{\raggedright}p{5.2cm} >{\raggedright}p{1.8cm} >{\raggedright\arraybackslash}p{7.6cm}}
\toprule \textbf{Entregable} & \textbf{Estado} & \textbf{Evidencia / observación}\\ \midrule \endhead}
{\bottomrule\end{longtable}\end{center}}

\input{../../gemelo-digital/manual/estilo_manual}
% Macros del manual técnico: ladder, estados, código y rutas de imágenes
\usetikzlibrary{shapes.multipart}
\graphicspath{{../img/}{../../../../../global/fotos-planta/}{../../../modelo-3d/informe/figuras/}{../../../modelo-3d/levantamiento/}{../../../database/modelo/}{../../../documentacion/manuales_inversa/}{../../../documentacion/planos_isa/electrico/}{../../../database/interfaz/}}

\lstdefinestyle{py}{language=Python, keywordstyle=\color{azulg}\bfseries, commentstyle=\color{grisos}\itshape,
  stringstyle=\color{verde}, basicstyle=\ttfamily\scriptsize, numbers=left, numberstyle=\tiny\color{grisos}}
\lstdefinestyle{js}{language=Java, morekeywords={const,let,function,async,await,return,for,of,if,else}, keywordstyle=\color{azulg}\bfseries,
  commentstyle=\color{grisos}\itshape, stringstyle=\color{verde}, basicstyle=\ttfamily\scriptsize, numbers=left, numberstyle=\tiny\color{grisos}}
\lstdefinestyle{awl}{basicstyle=\ttfamily\tiny, commentstyle=\color{grisos}\itshape, morecomment=[l]{//},
  keywordstyle=\color{azulg}\bfseries, morekeywords={LD,LDN,A,AN,O,ON,EU,ED,S,R,MOVB,MOVW,MOVD,PLS,TON,CALL,CRET,LPS,LRD,LPP,ALD,OLD,FILL,LDW,LDB,LDD,AW,AD,AB,OW,OD},
  numbers=left, numberstyle=\tiny\color{grisos}, frame=single, backgroundcolor=\color{gris}}
\lstdefinestyle{sql}{language=SQL, keywordstyle=\color{azulg}\bfseries, commentstyle=\color{grisos}\itshape,
  stringstyle=\color{verde}, basicstyle=\ttfamily\tiny, numbers=left, numberstyle=\tiny\color{grisos},
  morekeywords={REFERENCES,SERIAL,SMALLSERIAL,BIGSERIAL,TIMESTAMPTZ,BOOLEAN,INET,RETURNS,TRIGGER,LANGUAGE,DECLARE,RAISE,EXCEPTION,VIEW,INDEX}}

\tikzset{
  ld/.style={line width=0.7pt},
  pics/na/.style={code={\fill[white](-0.12,-0.26) rectangle (0.12,0.26);
     \draw[thick](-0.1,-0.24)--(-0.1,0.24)(0.1,-0.24)--(0.1,0.24);
     \node[font=\scriptsize\ttfamily, above] at (0,0.22){#1};}},
  pics/nc/.style={code={\fill[white](-0.12,-0.26) rectangle (0.12,0.26);
     \draw[thick](-0.1,-0.24)--(-0.1,0.24)(0.1,-0.24)--(0.1,0.24);\draw(-0.17,-0.24)--(0.17,0.24);
     \node[font=\scriptsize\ttfamily, above] at (0,0.22){#1};}},
  pics/pf/.style={code={\fill[white](-0.12,-0.26) rectangle (0.12,0.26);
     \draw[thick](-0.1,-0.24)--(-0.1,0.24)(0.1,-0.24)--(0.1,0.24);
     \node[font=\tiny\bfseries] at (0,0){P};}},
  pics/cmp/.style n args={3}{code={\fill[white](-0.12,-0.26) rectangle (0.12,0.26);
     \draw[thick](-0.1,-0.24)--(-0.1,0.24)(0.1,-0.24)--(0.1,0.24);
     \node[font=\tiny\ttfamily] at (0,0){#2};
     \node[font=\scriptsize\ttfamily, above] at (0,0.22){#1};
     \node[font=\scriptsize\ttfamily, below] at (0,-0.22){#3};}},
  pics/bob/.style n args={3}{code={\fill[white](-0.2,-0.26) rectangle (0.2,0.26);
     \draw[thick](-0.05,0.24) arc (110:250:0.255);\draw[thick](0.05,0.24) arc (70:-70:0.255);
     \node[font=\tiny\bfseries] at (0,0){#2};
     \node[font=\scriptsize\ttfamily, above] at (0,0.22){#1};
     \node[font=\scriptsize\ttfamily, below] at (0,-0.22){#3};}},
  fb/.style={draw, fill=white, font=\scriptsize\ttfamily, align=left, inner sep=3pt, anchor=west},
  barra/.style={line width=2pt},
  estado/.style={draw, rounded corners=3pt, fill=white, align=center, font=\scriptsize, minimum width=22mm, minimum height=10mm, inner sep=3pt},
  estadon/.style={estado, fill=naranja!14, draw=naranja!70!black},
  trans/.style={font=\scriptsize, align=center, fill=white, inner sep=1pt},
  lugar/.style={circle, draw, thick, minimum size=7mm, fill=white},
  transicion/.style={rectangle, fill=black, minimum width=1.6mm, minimum height=7mm, inner sep=0},
  tabla/.style={rectangle split, rectangle split parts=2, draw, fill=white, font=\tiny\ttfamily, align=left, inner sep=2pt, rounded corners=1pt, text width=33mm},
}
\newcommand{\pendiente}[1]{\par\smallskip\noindent\colorbox{naranja!12}{\parbox{\dimexpr\linewidth-2\fboxsep}{\textbf{\color{naranja!80!black}Pendiente.} #1}}\par\smallskip}
\newcommand{\cumple}{\textcolor{verde}{\textbf{cumplido}}}
\newcommand{\parcial}{\textcolor{naranja}{\textbf{parcial}}}
\newcommand{\falta}{\textcolor{rojo}{\textbf{pendiente}}}

\graphicspath{{../../documentacion/planos_isa/pid/}{../../documentacion/planos_isa/electrico/}{../../documentacion/manuales_inversa/}
  {../../database/mockups/}{../../database/documentacion/}{../../database/modelo/}{../../database/interfaz/}
  {../../gemelo-digital/manual/img/}{../../modelo-3d/informe/figuras/}{../../modelo-3d/levantamiento/}{../../../../global/fotos-planta/}{../figuras/}}
\usepackage{pdflscape}
\lstset{showstringspaces=false}
% plano A3 en una página apaisada: \plano[página]{archivo}{leyenda}
\newcommand{\plano}[3][1]{\begin{landscape}\begin{figure}[p]\centering\includegraphics[page=#1, width=\linewidth, height=0.86\textwidth, keepaspectratio]{#2}\caption{#3}\end{figure}\end{landscape}}
\setlength{\headheight}{14pt}
\newcommand{\nlab}{}
\newcommand{\tlab}{}
\lhead{\small Brazo XK-JS1-1 · Informe del laboratorio \nlab}
\rhead{\small UMSA · II/2026}

% \portada{número}{título}{inicio}{entrega}{semanas}
\newcommand{\portada}[5]{%
\renewcommand{\nlab}{#1}%
\begin{titlepage}
\centering
\vspace*{0.6cm}
{\large Universidad Mayor de San Andrés · Facultad de Ingeniería\par}
{\large Carrera de Ingeniería Electrónica\par}
\vspace{0.8cm}
{\color{naranja!80!black}\rule{\linewidth}{1.5pt}\par}
\vspace{0.3cm}
{\Large Informe del Laboratorio Integrado Nº #1\par}
\vspace{0.2cm}
{\huge\bfseries #2\par}
\vspace{0.3cm}
{\color{naranja!80!black}\rule{\linewidth}{1.5pt}\par}
\vspace{0.8cm}
{\large Planta: brazo robótico XK-JS1-1 de 4 grados de libertad\par}
\vspace{0.3cm}
{\large ETN-902 Sistemas de Control II · ETN-1034 Aplicación de Técnicas de Control · ETN-1000 Bases de Datos\par}
\vspace{1cm}
\begin{tabular}{rl}
\textbf{Docente:} & Ing. Jorge Antonio Nava Amador\\
\textbf{Auxiliares:} & Limber Ajnota · Fabián Plata · Franz Choque\\
\textbf{Inicio de trabajo:} & #3\\
\textbf{Entrega:} & #4 (#5)\\
\end{tabular}
\vspace{1cm}

\begin{tabular}{ll}
\multicolumn{2}{l}{\textbf{Escuadra del brazo de 4 GDL}}\\[2pt]
Control (ETN-902 / ETN-1034): & Gabriel Andres Gomez Calla\\
 & Jose Lenin Mamani Choque\\
 & Miguel Alejandro Lima Huarachi\\
 & Román Abner Machicado Arequipa\\[2pt]
Cómputo (ETN-1000): & Miguel Alejandro Lima Huarachi\\
 & Marco Eduardo Chura Rondo\\
\end{tabular}
\vfill
{\small Repositorio: \url{https://github.com/GabAndGom/Brazo-de-4-grados-de-libertad}\par}
{\small La Paz, Bolivia · gestión II/2026\par}
\end{titlepage}
}

% Tabla de entregables (Definition of Done): \begin{dod} ... \end{dod}, filas: entregable & estado & evidencia\\
\newenvironment{dod}{\begin{center}\small\begin{longtable}{>{\raggedright}p{5.2cm} >{\raggedright}p{1.8cm} >{\raggedright\arraybackslash}p{7.6cm}}
\toprule \textbf{Entregable} & \textbf{Estado} & \textbf{Evidencia / observación}\\ \midrule \endhead}
{\bottomrule\end{longtable}\end{center}}

\renewcommand{\tipodoc}{Logros del Laboratorio Integrado}
\renewcommand{\tipocab}{Logros del laboratorio}
\begin{document}
\portada{3}{Control de versiones con Git y modelo entidad-relación}{martes 1 de septiembre de 2026}{martes 15 de septiembre de 2026}{semanas 5 y 6}
\tableofcontents
\newpage
\nota{\textbf{Estado al 18/11/2026.} Laboratorio entregado; cierra la fase 1. Fase 1: ingeniería inversa y auditoría Kaizen (cierre). Este documento explica, objetivo por objetivo, qué pide la guía, qué logró la escuadra, qué es cada cosa y dónde está la evidencia; después explica el trabajo del docente: qué debía revisar en cada hito y qué tiene preparado la escuadra para esa revisión.}
\section{Resumen}
\begin{center}\small\begin{tabularx}{\linewidth}{c >{\raggedright}p{2.4cm} >{\raggedright}X c c}
\toprule \textbf{Nº} & \textbf{Mención} & \textbf{Objetivo} & \textbf{Estado} & \textbf{\%}\\ \midrule
1 & ETN-902 / ETN-1034 & Gestión del ciclo de vida del software de control & \parcial & 80\%\\
2 & ETN-902 / ETN-1034 & Gestión de cambios de ingeniería (Pull Requests) & \parcial & 40\%\\
3 & ETN-1000 & Modelado relacional de datos & \cumple & 100\%\\
4 & ETN-1000 & Estandarización de repositorios & \cumple & 100\%\\
\bottomrule\end{tabularx}\end{center}
El porcentaje mide cuánto del objetivo está hecho y comprobado en la planta o en el repositorio; no cuenta lo que solo está planificado.
\section{Logros respecto de los objetivos}
\subsection{Mención Control Automático (ETN-902 / ETN-1034)}
\subsubsection{Gestión del ciclo de vida del software de control}
\paragraph{Qué pide la guía.} Custodiar la documentación y los scripts con Git y GitHub.
\paragraph{Qué se logró.}
\begin{itemize}
\item Repositorio público en GitHub con 63 commits en \codigo{main}; cada cambio de planos, PLC, gemelo y base de datos queda registrado con su mensaje.
\item Historial generado con \codigo{git log --oneline --graph} en \codigo{git\_log.txt}.
\item La página del proyecto (brazo-4gdl.pages.dev) se publica sola con cada cambio en \codigo{main}.
\end{itemize}
\paragraph{Qué es y por qué importa.} Git guarda cada versión de cada archivo con quién lo cambió, cuándo y por qué. Con eso no se pierde trabajo entre semestres y se puede volver a cualquier versión anterior de un plano o programa.
\paragraph{Evidencia.}
\begin{itemize}\small
\item \url{https://github.com/GabAndGom/Brazo-de-4-grados-de-libertad}
\item \codigo{laboratorios/lab03\_git\_y\_modelo\_er/git\_log.txt}
\end{itemize}
\paragraph{Lo que falta.} Actividad de todos los integrantes: hoy Gabriel 58, Jose Lenin 2, Román 2, Eduardo 1 y Miguel 0 commits.
\paragraph{Estado:} \parcial~(80\%).
\subsubsection{Gestión de cambios de ingeniería (Pull Requests)}
\paragraph{Qué pide la guía.} Consolidar la v2.0 de los planos con un flujo formal de revisión y aprobación (Pull Requests).
\paragraph{Qué se logró.}
\begin{itemize}
\item Planos v2.0 consolidados en \codigo{documentacion/planos\_isa/}.
\item 4 ramas por componente: \codigo{feature/planos-v2-kaizen}, \codigo{feature/ladder-s7200-homing}, \codigo{feature/db-y-supervision} y \codigo{feat/cinematica-modelo-3d}; una fusión registrada (\codigo{ac8ab17}).
\end{itemize}
\paragraph{Qué es y por qué importa.} Un Pull Request es una solicitud de fusionar una rama en \codigo{main} que otra persona revisa antes de aceptarla. Simula la oficina de ingeniería: nadie cambia el plano oficial sin que un revisor verifique que cumple ISA 5.1 e IEC.
\paragraph{Evidencia.}
\begin{itemize}\small
\item \codigo{laboratorios/lab03\_git\_y\_modelo\_er/informe\_lab03.pdf}
\end{itemize}
\paragraph{Lo que falta.} Abrir y fusionar al menos un Pull Request por rama con el líder o el auxiliar como revisor (hoy hay 0).
\paragraph{Estado:} \parcial~(40\%).
\subsection{Mención Sistemas Computacionales (ETN-1000)}
\subsubsection{Modelado relacional de datos}
\paragraph{Qué pide la guía.} Convertir los requerimientos y mockups en un modelo E-R y lógico normalizado (1FN, 2FN, 3FN).
\paragraph{Qué se logró.}
\begin{itemize}
\item Modelo E-R y lógico de 15 tablas en 3FN: catálogo de la planta (sistema, eje, punto\_io, variable), operación (ciclo\_operacion, orden, orden\_parametro), telemetría (muestra, estado\_io), alarmas (alarma, evento\_alarma) y usuarios (usuario, sesion).
\item Cada tabla con PK, FK, tipos y cardinalidad (1:N y la M:N orden $\leftrightarrow$ parámetro resuelta con tabla intermedia).
\item Justificación de cada forma normal en el documento de base de datos (25 páginas) y diccionario de datos generado del script.
\end{itemize}
\paragraph{Qué es y por qué importa.} Normalizar es repartir los datos en tablas de modo que cada dato se guarde una sola vez: 1FN (un valor por celda), 2FN (todo depende de la clave completa) y 3FN (nada depende de otra columna que no sea la clave). Así no hay duplicados que se contradigan.
\paragraph{Evidencia.}
\begin{itemize}\small
\item \codigo{database/documentacion/base\_datos\_brazo.pdf}
\item \codigo{database/scripts/brazo\_4gdl.sql}
\end{itemize}
\paragraph{Lo que falta.} nada para este objetivo.
\paragraph{Estado:} \cumple~(100\%).
\subsubsection{Estandarización de repositorios}
\paragraph{Qué pide la guía.} Estructura de carpetas modular donde convivan base de datos, middleware y archivos CAD.
\paragraph{Qué se logró.}
\begin{itemize}
\item Estructura de la guía: \codigo{documentacion/planos\_isa/}, \codigo{documentacion/manuales\_inversa/}, \codigo{database/scripts/}, más una carpeta por laboratorio, \codigo{gemelo-digital/}, \codigo{programas-plc/} y \codigo{modelo-3d/}.
\item README en cada carpeta explicando su contenido; la gestión 2026\_I se conserva sin cambios.
\end{itemize}
\paragraph{Qué es y por qué importa.} Una estructura acordada hace que cualquier integrante (o el grupo del próximo semestre) encuentre cada archivo sin preguntar, y evita que el código de la base, el middleware y los planos se mezclen.
\paragraph{Evidencia.}
\begin{itemize}\small
\item \codigo{semestres/2026\_II/}
\item \codigo{laboratorios/README.md}
\end{itemize}
\paragraph{Lo que falta.} nada para este objetivo.
\paragraph{Estado:} \cumple~(100\%).
\section{Trabajo del docente}
La guía asigna al docente el papel de auditor: en cada hito revisa el trabajo de la escuadra y decide si cumple antes de pasar a la etapa siguiente. Para cada tarea se explica qué debía revisar, cuándo, y qué tiene la escuadra para presentarle. El porcentaje es el del Excel \codigo{evaluacion\_tareas\_docente.xlsx}.
\subsection{Auditoría de repositorios}
\paragraph{Lo que dice la guía.} Inspeccionar los perfiles de GitHub el martes de la semana 6: frecuencia, claridad y pertinencia de los commits. \textit{Cuándo:} martes 8 de septiembre de 2026 (semana 6).
\paragraph{Qué debía revisar el docente (hitos de la revisión).}
\begin{itemize}
\item[$\square$] Que todos los integrantes tengan commits propios (la guía lo exige sin excepción).
\item[$\square$] Que los mensajes de commit sean descriptivos («Planos v2.0: P\&ID ISA-5.1…») y no «cambios» o «update».
\item[$\square$] Que existan ramas por componente y Pull Requests revisados antes de fusionar (30 \% de la defensa).
\end{itemize}
\paragraph{Qué tiene la escuadra para esa revisión.}
\begin{itemize}
\item 4 ramas por componente y \codigo{git\_log.txt}.
\item Commits en main: Gabriel 58, Jose Lenin 2, Román 2, Eduardo 1, Miguel 0. Pull Requests: 0.
\end{itemize}
\paragraph{Estado:} \parcial~(50\%).
\subsection{Filtro de la fase 1 (E-R frente a tags ISA)}
\paragraph{Lo que dice la guía.} Evaluar que el modelo E-R refleje fielmente los tags ISA 5.1 de los lazos de control. \textit{Cuándo:} martes 15 de septiembre de 2026 (cierre de fase 1).
\paragraph{Qué debía revisar el docente (hitos de la revisión).}
\begin{itemize}
\item[$\square$] Que cada señal del P\&ID (ZS-101, ZSL-102, ZY-102…) exista con el mismo tag en la base.
\item[$\square$] Que las variables PV y MV tengan su lazo y unidad, y que el modelo esté en 3FN (40 \% de la defensa).
\end{itemize}
\paragraph{Qué tiene la escuadra para esa revisión.}
\begin{itemize}
\item 22 tags de señales en \codigo{punto\_io} y 5 de variables en \codigo{variable}, únicos, con su lazo en \codigo{eje.lazo\_isa}.
\item Tabla de correspondencia plano $\leftrightarrow$ base en el informe 3.
\end{itemize}
\paragraph{Estado:} \cumple~(100\%).
\subsection{Repositorios y control de versiones inicial (tarea de la fase 1)}
\paragraph{Lo que dice la guía.} Establecer los repositorios locales y el control de versiones inicial. \textit{Cuándo:} fase 1 (semanas 1 a 6).
\paragraph{Qué debía revisar el docente (hitos de la revisión).}
\begin{itemize}
\item[$\square$] Que cada escuadra tenga su repositorio con la estructura de la guía 3 y acceso de todos los integrantes.
\end{itemize}
\paragraph{Qué tiene la escuadra para esa revisión.}
\begin{itemize}
\item Repositorio público con 63 commits, 4 ramas y una carpeta por laboratorio.
\end{itemize}
\paragraph{Estado:} \parcial~(90\%).
\subsection{Hito de la fase}
\textbf{Hito de la fase 1 (fin de la semana 6):} versión 2.0 corregida de los P\&ID y esquemas eléctricos. Para el brazo está cumplido: P\&ID y eléctrico v2.0 en DWG, DXF y PDF, con las correcciones verificadas. El informe de cierre de fase pide además los perfiles ORCID del equipo, que todavía faltan.
\section{Conclusión}
Cumplidos: Modelado relacional de datos, Estandarización de repositorios. Con avance parcial: Gestión del ciclo de vida del software de control, Gestión de cambios de ingeniería (Pull Requests).
\par\smallskip Pendientes principales:
\begin{enumerate}
\item Actividad de todos los integrantes: hoy Gabriel 58, Jose Lenin 2, Román 2, Eduardo 1 y Miguel 0 commits.
\item Abrir y fusionar al menos un Pull Request por rama con el líder o el auxiliar como revisor (hoy hay 0).
\end{enumerate}
\par\smallskip\noindent Repositorio: \url{https://github.com/GabAndGom/Brazo-de-4-grados-de-libertad}
\end{document}
